\chapter{Catalogues finis et certificats de reproductibilité}
\label{app:datos-reproducibles}

Les objets finis suivants permettent de répéter les vérifications sans
incorporer la notation de production à l'exposé mathématique. Chaque contrôle
conserve le domaine, l'opération, l'invariant et le réfutateur qui lui
correspondent.

\section{Catalogues finis de l'état commun}

\begin{longtable}{@{}>{\raggedright\arraybackslash}p{.34\textwidth}
  >{\raggedright\arraybackslash}p{.58\textwidth}@{}}
\toprule
Objet fini & Contenu mathématique vérifiable\\
\midrule
\endhead
Catalogue des émissions de longueur six
& Les \(468\) émissions \(U_6\), leurs multiplicités, leur orientation et leurs signatures
additive et multiplicative.\\
Partition spectrale des émissions
& Les vingt-six classes de taille \(18\), avec l'application explicite qui
reconstruit chaque émission à partir de sa classe et de sa position interne.\\
Atlas dodécaphasique des voies
& Les \(468\) voies minimales avec feuillet, phase, retenue, frontière et mémoire.\\
Catalogue du lecteur archimédien
& Les cinq régions de \(\pi\), les régions conjuguées de \(e\) et l'axe
d'auto-échelle de \(\varphi\), avec leurs cylindres et multiplicités.\\
Flux de Hensel de vingt blocs
& La matrice unimodulaire, son inverse entier, les états \(3\)-adiques et la
récurrence reste--quotient qui prolonge les trois canaux.\\
Registre de monodromie nonadique
& Les vingt blocs, leurs phases, préfixes ternaires et préfixes décimaux ;
il inclut les \(43\) retours complets et les \(387\) paires \((K,K+9)\) par
canal.\\
Fenêtre de survie \(K=5,\ldots,13\)
& Les candidats locaux, la branche compatible et les trois sorties de la
transition \(9\to1\), sans identifier le retour de phase à la réinitialisation de l'état.\\
Lecture archimédienne initiale
& Les douze triplets décimaux obtenus à partir des treize premiers blocs de chaque
canal, séparés du certificat de précision arbitraire.\\
\bottomrule
\end{longtable}

\section{Identités algébriques, opératorielles et spectrales des certificats finis}

\begin{longtable}{@{}>{\raggedright\arraybackslash}p{.34\textwidth}
  >{\raggedright\arraybackslash}p{.58\textwidth}@{}}
\toprule
Contrôle & Identités et réfutateurs\\
\midrule
\endhead
Monodromie et génération des constantes
& Vérifie \(w_6\to w_{30}\to R_{36}\to G_9\), la sélection \(3\to1\),
le changement d'état après neuf phases, les cinq régions de \(\pi\) et la
compatibilité des troncatures. Échoue si un chiffre cible est introduit
avant le lecteur.\\
Coordonnée angulaire conjuguée
& Vérifie le sélecteur régional \(169\), la récurrence déterminantale, la
somme analytique et sa borne de reste, et l'unicité de \(C^*\). Échoue si une
constante du SI intervient dans la sélection.\\
Dynamique du quotient \(468=18\cdot26\)
& Certifie la bijection entre émissions et classes, le retour uniforme, la
mesure de l'ombre moyennée, les cocycles de mode et le relèvement
microscopique. Distingue la fréquence chronologique de la mesure de Parry.\\
Canaux entiers des constantes
& Vérifie les identités \((729,271,315,161)\), les réalisations
ensemblistes de \(271/315\), les orbites diédriques et la factorisation de la
parité visible et de la parité de retenue.\\
Incidence de Witt
& Reconstruit l'incidence point--hexade, le module \(E_{30}\), le
transport positif de \(P_3\), l'étoile de Witt et les obstructions
d'équivariance.\\
Sélecteur variationnel
& Reconstruit le normalisateur \(S_3\), la moyenne de Reynolds, la
factorisation polaire et la chaîne de projecteurs de rang trois ; la positivité
est prouvée exactement dans \(\mathbb Q(\sqrt5)\).\\
Relèvement \(W_{12}\to W_{24}\)
& Vérifie les incidences de quatre et cinq points, le relèvement unique des
hexades et la correspondance typée avec les cinq régions de \(\pi\).\\
Realizaciones de \(-1/12\)
& Construit les projecteurs rationnels des décalages de pas trois
et quatre, vérifie la symétrie et l'idempotence et obtient
\((3-4)/12=-1/12\) comme différence de traces normalisées.\\
Holonomie de Schläfli
& Parcourt les plaquettes, certifie \(A^2=5I-4A+5J\), la configuration des
vingt-sept droites, leurs \(45\) triangles et la partition
\(729=27+270+432\).\\
Fonctionnelles spectrales
& Recalcule Catalan, Stieltjes, Apéry, Bendersky--Glaisher--Kinkelin, les
deux extensions d'Euler--Kronecker et les bornes dirigées des produits premiers.\\
Lacunes et criticité
& Vérifie les séparations \(21/22\), les retours \(6/7\), la partie
harmonique finie et les approximants dodécaphasiques du mode critique.\\
\bottomrule
\end{longtable}

\section{Dictionnaire typé des catalogues}

\begin{longtable}{>{\ttfamily}p{.18\textwidth}>{\raggedright\arraybackslash}p{.72\textwidth}}
\toprule
Champ & Signification et codomaine\\
\midrule
\endhead
Esig & Signature du canal multiplicatif ; identifie l'une des vingt-six
classes de la factorisation de l'émetteur.\\
hplus\_type & Classe d'orientation du canal additif pour les deux paires de
directions cardinales déclarées.\\
diag\_c & Indice entier de la classe diagonale du curseur additif.\\
px\_size & Cardinal de la fibre de la classe additive correspondante.\\
colw & Sextuplet des poids de colonne du bloc émis ; appartient à
\(\mathbb Z_{\geq0}^{6}\).\\
\bottomrule
\end{longtable}

\section{Reproductibilité des certificats et obstructions exactes des constructions finies}

Les contrôles utilisent l'arithmétique entière ou rationnelle chaque fois que l'objet le
permet. Les bornes analytiques utilisent des intervalles avec arrondi dirigé et séparent
la démonstration symbolique de la comparaison décimale. Chaque exécution
reconstruit les objets à partir des germes déclarés, vérifie les identités
d'incidence et de naturalité, et confronte les sorties à une implémentation
arithmétique indépendante.

Le contrôle du produit natif rejette toute dépendance envers une évaluation
décimale, une factorisation externe, un crible, une table de zéros ou une
constante cible. Le contrôle nonadique rejette l'identification de la fibre
ambiante de \(729\) éléments au cardinal d'un espace d'histoires. Le
contrôle de phase rejette \(g(K+9)=g(K)\) comme preuve d'égalité des états. Le
contrôle métrologique rejette toute sélection de voie effectuée après
consultation d'une grandeur observée.

Une divergence entre une identité du texte, le catalogue fini et son contrôle
indépendant réfute la certification correspondante. Cette règle ne remplace
aucune démonstration : elle fixe un second canal de vérification pour les objets
finis qui y interviennent.
